% ============================================================
%  chuleta_hibridacion_enlaces.tex
%  Chuleta de Química: Hibridación y Enlaces
%  Mathwizards Consultoría Educativa STEM
%
%  Compilar:  latexmk -xelatex chuleta_hibridacion_enlaces.tex
% ============================================================
\documentclass[9pt, landscape, letterpaper]{extarticle}
\usepackage{mathwizards-formulario}

\mwformulario{Chuleta de Química: Hibridación y Enlaces}{Química $\cdot$ Enlace Químico}

\begin{document}

% ════════════════════════════════════════════════════════════
%  PÁGINA 1 — Hibridación y Tipos de Enlace
% ════════════════════════════════════════════════════════════
\nuevotema{Hibridación y Tipos de Enlace}

\begin{multicols}{2}

  \subsection{¿Qué es la hibridación?}

  \begin{notabox}[Idea clave]
    Es la \textbf{mezcla de orbitales atómicos} para formar orbitales
    híbridos \textbf{equivalentes}. Explica la geometría real de la
    molécula y el ángulo entre sus enlaces.
  \end{notabox}

  % ── Hibridación sp ────────────────────────────────────────
  \subsection{Hibridación \textbf{sp} — Lineal}

  \begin{formulabox}[Ángulo: $180^\circ$]
    \centering
    \begin{tikzpicture}[scale=0.62, baseline=(current bounding box.center)]
      \draw[mwprimario, line width=1.2pt] (-2.0,0) -- (2.0,0);
      \fill[mwprimario] (0,0) circle (2.8pt);
      \node[mwprimario, font=\scriptsize] at (0,0.62) {$180^\circ$};
    \end{tikzpicture}
    \\[2pt]
    {\small Se forma con \textbf{dos dobles enlaces} \textbf{o} con \textbf{un triple enlace}.}
  \end{formulabox}

  % ── Hibridación sp2 ───────────────────────────────────────
  \subsection{Hibridación \textbf{sp\textsuperscript{2}} — Triangular Plana}

  \begin{formulabox}[Ángulo: $120^\circ$]
    \centering
    \begin{tikzpicture}[scale=0.62, baseline=(current bounding box.center)]
      \foreach \a in {90, 210, 330}{%
        \draw[mwprimario, line width=1.2pt] (0,0) -- (\a:1.9);}
      \fill[mwprimario] (0,0) circle (2.8pt);
      \node[mwprimario, font=\scriptsize, anchor=west] at (0.12,0.55) {$120^\circ$};
    \end{tikzpicture}
    \\[2pt]
    {\small Se forma con \textbf{un solo enlace doble} (el resto, simples).}
  \end{formulabox}

  % ── Hibridación sp3 ───────────────────────────────────────
  \subsection{Hibridación \textbf{sp\textsuperscript{3}} — Tetraédrica}

  \begin{formulabox}[Ángulo: $109{,}5^\circ$]
    \centering
    \begin{tikzpicture}[scale=0.62, baseline=(current bounding box.center)]
      \draw[mwprimario, line width=1.2pt] (0,0) -- (90:1.8);
      \draw[mwprimario, line width=1.2pt] (0,0) -- (210:1.8);
      \draw[mwprimario, line width=1.2pt] (0,0) -- (330:1.8);
      \draw[mwprimario, line width=1.0pt, dashed] (0,0) -- (48:1.05);
      \fill[mwprimario] (0,0) circle (2.8pt);
      \node[mwprimario, font=\scriptsize, anchor=west] at (0.12,0.18) {$109{,}5^\circ$};
    \end{tikzpicture}
    \\[2pt]
    {\small Se forma cuando \textbf{todos los enlaces son simples}.}
  \end{formulabox}

  % ── Tabla resumen ─────────────────────────────────────────
  \subsection{Resumen de Hibridaciones}

  \begin{tablabox}[Geometría según hibridación]
    \tablacompacta
    \begin{tabular}{@{}l l c l@{}}
      \toprule
      Hibrid. & Geometría        & Ángulo           & Se forma con        \\
      \midrule
      $sp$    & Lineal           & $180^\circ$      & 2 dobles ó 1 triple \\
      $sp^2$  & Triangular plana & $120^\circ$      & 1 enlace doble      \\
      $sp^3$  & Tetraédrica      & $109{,}5^\circ$  & Todos simples       \\
      \bottomrule
    \end{tabular}
  \end{tablabox}

  % ── Tipos de enlace ───────────────────────────────────────
  \subsection{Composición del Enlace}

  \begin{tablabox}[Tipo de enlace]
    \tablacompacta
    \begin{tabular}{@{}l l@{}}
      \toprule
      Tipo de enlace & Composición     \\
      \midrule
      \textbf{Simple}  & 1 $\sigma$      \\
      \textbf{Doble}   & 1 $\sigma$ + 1 $\pi$ \\
      \textbf{Triple}  & 1 $\sigma$ + 2 $\pi$ \\
      \bottomrule
    \end{tabular}
  \end{tablabox}

  \begin{notabox}[Recuerda]
    El \textbf{primer} enlace entre dos átomos siempre es \textbf{sigma}
    ($\sigma$); los enlaces \textbf{adicionales} son \textbf{pi} ($\pi$).
  \end{notabox}

  % ── Ejemplos por hibridación ──────────────────────────────
  \subsection{Ejemplos Típicos}

  \begin{tablabox}[Moléculas según hibridación]
    \tablacompacta
    \begin{tabular}{@{}l l@{}}
      \toprule
      Hibrid. & Moléculas                      \\
      \midrule
      $sp$    & CO$_2$, C$_2$H$_2$ (etino)     \\
      $sp^2$  & C$_2$H$_4$ (eteno), benceno    \\
      $sp^3$  & CH$_4$, C$_2$H$_6$, H$_2$O, NH$_3$ \\
      \bottomrule
    \end{tabular}
  \end{tablabox}

  \subsection{¿Cómo determinarla?}

  \begin{notabox}[Cuenta los dominios electrónicos]
    Suma los \textbf{enlaces} ($\sigma$) más los \textbf{pares libres}
    del átomo central:
    \begin{itemize}
      \item \textbf{2 dominios} $\to$ $sp$ (lineal).
      \item \textbf{3 dominios} $\to$ $sp^2$ (triangular plana).
      \item \textbf{4 dominios} $\to$ $sp^3$ (tetraédrica).
    \end{itemize}
  \end{notabox}

\end{multicols}

% ════════════════════════════════════════════════════════════
%  PÁGINA 2 — Diferencia de Electronegatividad (ΔEN)
% ════════════════════════════════════════════════════════════
\nuevotema{Diferencia de Electronegatividad ($\Delta EN$)}

\begin{multicols}{2}

  \subsection{Tipo de Enlace según $\Delta EN$}

  \begin{tablabox}[$\Delta EN = |EN_1 - EN_2|$]
    \tablacompacta
    \begin{tabular}{@{}l l l@{}}
      \toprule
      $\Delta EN$           & Tipo de enlace              & Ejemplo              \\
      \midrule
      $0$                   & Covalente no polar (puro)   & H$_2$, O$_2$, N$_2$, Cl$_2$ \\
      $0{,}1 - 0{,}4$       & Covalente no polar          & C--H ($0{,}4$)       \\
      $0{,}5 - 1{,}7$       & Covalente polar             & H--Cl ($0{,}9$), H$_2$O \\
      $> 1{,}7$             & Iónico                      & NaCl ($2{,}1$), MgO  \\
      $\approx 0$ (metal--metal) & Metálico               & Fe, Cu, Au, aleaciones \\
      \bottomrule
    \end{tabular}
  \end{tablabox}

  \subsection{Características de cada Enlace}

  \begin{notabox}[¿Cómo se comportan los electrones?]
    \begin{itemize}
      \item \textbf{No polar:} los electrones se comparten
            \emph{equitativamente}.
      \item \textbf{Polar:} se comparten de forma \emph{desigual}
            (el más electronegativo atrae más).
      \item \textbf{Iónico:} hay \emph{transferencia} de electrones.
      \item \textbf{Metálico:} ``mar de electrones'' deslocalizados.
    \end{itemize}
  \end{notabox}

  \subsection{Escala de Pauling (Electronegatividades)}

  \begin{tablabox}[Valores comunes de $EN$]
    \tablacompacta
    \begin{tabular}{@{}l c l c@{}}
      \toprule
      Elemento & EN   & Elemento & EN   \\
      \midrule
      F  & 4{,}0 & C  & 2{,}5 \\
      O  & 3{,}5 & H  & 2{,}1 \\
      N  & 3{,}0 & Na & 0{,}9 \\
      Cl & 3{,}0 & Mg & 1{,}2 \\
      Br & 2{,}8 & K  & 0{,}8 \\
      I  & 2{,}5 & Ca & 1{,}0 \\
      S  & 2{,}5 & Fe & 1{,}8 \\
      P  & 2{,}1 & Cu & 1{,}9 \\
      \bottomrule
    \end{tabular}
  \end{tablabox}

  \subsection{Ejemplos de Cálculo}

  \begin{despejesbox}[NaCl — $\Delta EN = 3{,}0 - 0{,}9$]
    $$\Delta EN = \mathbf{2{,}1} \;\Rightarrow\; \text{\textbf{Enlace Iónico}}$$
  \end{despejesbox}

  \begin{despejesbox}[HCl — $\Delta EN = 3{,}0 - 2{,}1$]
    $$\Delta EN = \mathbf{0{,}9} \;\Rightarrow\; \text{\textbf{Covalente Polar}}$$
  \end{despejesbox}

  \begin{despejesbox}[O$_2$ — $\Delta EN = 3{,}5 - 3{,}5$]
    $$\Delta EN = \mathbf{0} \;\Rightarrow\; \text{\textbf{Covalente No Polar}}$$
  \end{despejesbox}

  \begin{despejesbox}[C--H — $\Delta EN = 2{,}5 - 2{,}1$]
    $$\Delta EN = \mathbf{0{,}4} \;\Rightarrow\; \text{\textbf{Covalente No Polar}}$$
  \end{despejesbox}

  \begin{despejesbox}[MgO — $\Delta EN = 3{,}5 - 1{,}2$]
    $$\Delta EN = \mathbf{2{,}3} \;\Rightarrow\; \text{\textbf{Enlace Iónico}}$$
  \end{despejesbox}

  \subsection{Regla Rápida}

  \begin{notabox}[Según los elementos]
    \begin{itemize}
      \item \textbf{No metal + No metal} $\to$ Covalente
            (polar o no polar según $\Delta EN$).
      \item \textbf{Metal + No metal} $\to$ Iónico
            (generalmente $\Delta EN > 1{,}7$).
      \item \textbf{Metal + Metal} $\to$ Metálico.
    \end{itemize}
  \end{notabox}

  \begin{notabox}[Enlace metálico]
    No se forma entre átomos aislados: los cationes metálicos quedan
    rodeados por un ``mar'' de electrones libres, lo que explica la
    \textbf{conductividad} y el \textbf{brillo} de los metales.
  \end{notabox}

\end{multicols}

\end{document}
